% Unit 32 — Series and complex exponentials
\input{technique-note.tex}
\title{Integration Bee --- Unit 32:\\ Series and Complex Exponentials}
\author{UK Uni Integration Bee prep\\ Termwise geometric series, and $\mathrm{Re}\,e^{(a+ib)x}$}
\date{}
\begin{document}
\maketitle
\section*{Status}
Tags: \texttt{yang:6.2}, \texttt{yang:6.3}; techniques: \texttt{taylor-series}, \texttt{complexifying}.
A geometric series on $[0,1)$ and Euler's formula $e^{ix}=\cos x+i\sin x$ are standard. Yang \S\S6.2--6.3 packages them as contest moves. Unit 08 is the first geometric integral. This unit says when the closed form is faster than the series, and when the complex exponential replaces two parts steps.
\section{Official anchors}
\begin{itemize}
\item Unit 08. The geometric sum $\sum_{n=0}^{N}(-1)^n x^n$ and its limit on $[0,1)$.
\item Syllabus \S1 series.
\item Yang Topics Guide 2nd Ed \S\S6.2 Taylor Series, 6.3 Complexifying.
\end{itemize}
\section{Learning goals}
\begin{enumerate}
  \item Integrate a geometric series term by term on $[0,1)$, then take the endpoint as a limit.
  \item If a closed antiderivative is one substitution, use it. The series explains the number. It is not always the fastest path.
  \item Replace $e^{ax}\cos bx$ by the real part of $e^{(a+ib)x}$.
  \item Refuse a series for a substitution, and refuse a definite value when the blank wants an antiderivative.
\end{enumerate}

\section{Worked series --- \texttt{U32-D01}}
\[
\int_0^1\frac{dx}{1+x}=\ln 2.
\]
The antiderivative is the fast path. The series path is the reason the value is the alternating harmonic number:
\[
\frac{1}{1+x}=\sum_{n=0}^{\infty}(-1)^n x^n,\qquad 0\le x<1,
\]
and $\int_0^1 x^n\,dx=1/(n+1)$, so the integral is $\sum_{n=0}^{\infty}(-1)^n/(n+1)=\ln 2$. At $x=1$ the geometric series diverges. On $[0,1]$ the partial sums are dominated by $2/(1+x)$, so term-by-term integration is valid and equals $\ln 2$. The worked dose stops there.

\section{Faded --- \texttt{U32-D02}}
\[
\int_0^1\Bigl(1-x^2+\frac{x^4}{2}-\frac{x^6}{6}\Bigr)dx.
\]
The blank is the four separate integrals. This is a polynomial. Refuse an infinite series. Integrate termwise:
\[
1-\frac13+\frac{1}{10}-\frac{1}{42}=\frac{26}{35}.
\]
That is the stored answer.

\section{Bare --- \texttt{U32-D04}}
Set
\[
\int_0^1 x^3\,dx
\]
and show nothing else. The bare dose hides the instruction to skip the series. After a line is written, uncover the stored answer $1/4$.

\section{Worked complex exponential --- \texttt{U32-D05}}
\[
\int_0^{\infty}e^{-x}\cos x\,dx.
\]
$e^{-x}\cos x$ is the real part of $e^{-x}e^{ix}=e^{-(1-i)x}$, and $e^{-x}\sin x$ is the imaginary part. The integral of the complex exponential is
\[
\int_0^{\infty}e^{-(1-i)x}\,dx=\frac{1}{1-i}=\frac{1+i}{2}.
\]
The real part is $1/2$ and the imaginary part is $1/2$. Those are the stored values of \texttt{U32-D05} and \texttt{U32-D06}. The same two numbers are \texttt{U30-D08} and \texttt{U30-D07}. The value $1/2$ does not record which part was taken. The indefinite pair below is where that choice is visible.

A larger decay is the fence on \texttt{docs/techniques/concepts/habits.md}, not a drill. For
\[
\int_0^{\infty}e^{-2x}\cos x\,dx,
\]
the complex integral is $1/(2-i)=(2+i)/5$, and the real part is $2/5$. SymPy checks $2/5$.

\section{Faded indefinite --- \texttt{U32-D07}}
\[
\int e^x\cos x\,dx.
\]
The blank is which part of $e^{(1+i)x}$ you keep. $e^x\cos x$ is the real part.
\[
\int e^{(1+i)x}\,dx=\frac{e^{(1+i)x}}{1+i}=\frac{1-i}{2}\,e^{(1+i)x}.
\]
Multiply $(1-i)(\cos x+i\sin x)$. The real part is $\cos x+\sin x$, so the stored antiderivative is
\[
\frac12 e^x(\cos x+\sin x)+C.
\]

\section{Bare --- \texttt{U32-D08}}
Set
\[
\int e^x\sin x\,dx
\]
and show nothing else. The bare dose hides the part. After a line is written, uncover the imaginary part of the same product. The stored antiderivative is
\[
\frac12 e^x(\sin x-\cos x)+C.
\]

\section{Refusal --- a series, or a definite number}
\texttt{U32-D03} is $\int x e^{x^2}\,dx$. One substitution stores $\frac12 e^{x^2}$. A series for $e^{x^2}$ is the long path. Refuse it.

\texttt{U32-D08} wants an antiderivative. Refuse the definite value $1/2$ from \texttt{U32-D05} or \texttt{U32-D06}. Refuse the cosine antiderivative of \texttt{U32-D07} when the integrand is sine.

\section{Dose}
For \texttt{taylor-series}, worked reads \texttt{U32-D01}, then the faded four terms on \texttt{U32-D02}, then sets \texttt{U32-D04}. Faded gives only those four integrals. Bare sets \texttt{U32-D04}. A series in place of a substitution opens the refusal and sets \texttt{U32-D03}.

For \texttt{complexifying}, worked reads \texttt{U32-D05}, then the faded part on \texttt{U32-D07}, then sets \texttt{U32-D08}. Faded gives only the choice of part. Bare sets \texttt{U32-D08}. A definite value, or the cosine antiderivative, written for the sine opens the refusal and sets \texttt{U32-D08}.

\section{Drill pointers}
See \texttt{drills.tex} / \texttt{drills.md}. Feed: \texttt{verify/drills\_unit32.json}.

\section*{On a miss}
Teach this file. Either write the geometric term and integrate it, or write $e^{(a+ib)x}$ and take the real or imaginary part. If a substitution already finished the integral, set \texttt{U32-D03}. If a definite value was written as an antiderivative, set \texttt{U32-D08}. Do not invent an answer that is not in \texttt{drills.md}.
\end{document}
